\chapter{Funcții de mai multe variabile}

\section{Limite și continuitate}

În cazul funcțiilor de mai multe variabile, atît studiul continuității,
cît și al derivabilității și calculului limitelor sau derivatelor
parțiale, se fac cu variabilele pe rînd, adică dînd prioritate
uneia dintre ele și considerînd pe celelalte drept parametri.

\subsection{Exerciții rezolvate}

1. Să se arate, folosind definiția, că următoarele funcții nu au limită în origine:
\begin{enumerate}[(a)]
\item $f(x, y) = \frac{2xy}{x^2 + y^2}, (x, y) \neq (0, 0)$;
\item $f(x, y) = \frac{x^2 + 2x}{y^2 - 2x}, (x, y) \neq (0, 0)$.
\end{enumerate}

\emph{Soluție:} Vom folosi definiția cu șiruri, punînd în evidență șiruri de puncte,
convergente către $(0, 0)$.

(a) Folosim șiruri de puncte care se află pe drepte care trec prin origine. Aceste drepte
au ecuații de forma $y = mx$. Dacă funcția ar avea limită în origine, ar trebui ca:
\[
  \lim_{(x, y) \to (0, 0)} f(x, y) = \lim_{x \to 0, y = mx} \frac{2x(mx)}{x^2 + x^2 m^2} = %
  \frac{2m}{1 + m^2}.
\]
Cu această limită depinde de șirul folosit (de $m$, mai exact), rezultă că limita nu
există, în general.

(b) Calculăm, similar:
\[
  \lim_{(x, y) \to (0, 0)} f(x, y) = %
  \lim_{x \to 0, y = mx} \frac{m^2 x^2 + 2x}{m^2 x^2 - 2x} = -1.
\]

Însă, dacă luăm $y^2 = px$, adică $(x, y) \to (0, 0)$ pe parabola care trece prin origine,
cu $p \in \RR - \{2 \}$, obținem:
\[
  \lim_{x \to 0, y^2 = px} \frac{px + 2x}{px - 2x} = \frac{p + 2}{p - 2}.
\]
Rezultă că, pentru șiruri diferite, obținem limite diferite, deci funcția nu are limită
în origine.

\vspace{1cm}

2. Să se calculeze limitele:
\begin{enumerate}[(a)]
\item $\ds\lim_{(x, y) \to (0, 0)} \frac{xy}{\sqrt{xy + 1} - 1}$;
\item $\ds\lim_{(x, y) \to (0, 2)} \frac{\sin xy}{x}$;
\item $\ds\lim_{(x, y) \to (\infty, \infty)} \frac{x + y}{x^2 + y^2}$.
\end{enumerate}

\emph{Soluție:} (a) Calculăm direct:
\begin{align*}
  \lim_{(x, y) \to (0, 0)} \frac{xy}{\sqrt{xy + 1} - 1} &= \lim_{u \to 0} \frac{u}{\sqrt{u + 1} - 1} \\
                                                        &= \lim_{u \to 0} \frac{(\sqrt{u + 1} + 1)u}{u}\\
                                                        &= 2.
\end{align*}

(b) Similar:
\[
  \lim_{(x, y) \to (0, 2)} \frac{\sin xy}{x} = \lim_{(x, y) \to (0, 2)} \frac{\sin xy}{xy} \cdot y = 2.
\]

(c) Pentru $x, y > 0$, avem:
\[
  0 < \frac{x + y}{x^2 + y^2} = \frac{x}{x^2 + y^2} + \frac{y}{x^2 + y^2} \leq \frac{x}{x^2} + \frac{y}{y^2}.
\]

Cum $\ds\lim_{(x, y) \to (\infty, \infty)} \Big( \frac{1}{x} + \frac{1}{y} \Big) = 0$, rezultă
că și limita inițială este nulă.

\vspace{1cm}

3. Să se arate că funcția:
\[
  f(x, y) = \begin{cases}
    \frac{xy^2 + \sin(x^3 + y^5)}{x^2 + y^4}, & (x, y) \neq (0, 0) \\
    0, & (x, y) = (0, 0)
  \end{cases}
\]
este continuă parțial în origine, dar nu este continuă în acest punct.

\emph{Soluție:} Considerăm variabilele separat. Funcția:
\[
  f(x, 0) = \begin{cases}
    \frac{\sin x^3}{x^2}, & x \neq 0 \\
    0, & x = 0
  \end{cases}
\]
este continuă în origine, deoarece $\ds\lim_{x \to 0} f(x, 0) = 0 = f(0, 0)$.

Similar, funcția $f(0, y)$ este continuă în $y = 0$.

Dar funcția $f(x, y)$ nu este continuă în origine, deoarece avem:
\begin{align*}
  \lim_{(x, y) \to (0, 0)} f(x, y) &= \lim_{x \to 0, y^2 = px} \frac{px^2 + \sin(x^3 + p^\frac{5}{2} x^{\frac{5}{2}})}{x^2 + p^2 x^2} \\
                                   &= \lim_{x \to 0, y^2 = px} \Bigg[ \frac{p}{1 + p^2} + %
                                     \frac{\sin(x^3 + p^{\frac{5}{2}} x^{\frac{5}{2}})}{x^3 + p^{\frac{5}{2}} x^{\frac{5}{2}}} \cdot %
                                     \frac{x^3 + p^{\frac{5}{2}} x^{\frac{5}{2}}}{x^2 + p^2 x^2} \Bigg] \\
                                   &= \frac{p}{1 + p^2},
\end{align*}
care depinde de $p$.

\vspace{1cm}

4. Să se arate că funcția:
\[
  f(x, y) = \begin{cases}
    \frac{2xy}{x^2 + y^2}, & (x, y) \neq (0, 0) \\
    0, & (x, y) = (0, 0)
  \end{cases}
\]
este continuă în raport cu fiecare variabilă în parte, dar nu este continuă în raport
cu ansamblul variabilelor.

\emph{Indicație:} Considerăm $y = b$, fixat. Avem $f(x, b)$ continuă. Similar, $f(a, y)$.

Pentru continuitatea globală în origine, luăm $y = mx$ și obținem o limită care depinde de $m$.

\vspace{1cm}

5. Să se arate că funcția:
\[
  f(x, y) = \begin{cases}
    \frac{xy + x^2 y \ln |x + y|}{x^2 + y^2}, & (x, y) \neq (0, 0) \\
    0, & (x, y) = (0, 0)
  \end{cases}
\]
este continuă parțial în origine, dar nu este continuă în raport cu ansamblul
variabilelor în acest punct.

\emph{Indicație:} Limitele iterate în origine sînt ambele nule, indiferent de ordine.

Însă pentru $x \to 0$ și $y = mx$, limita globală depinde de $m$.

%%%%%%%%%%%%%%%%%%%%%%%%%%%%%%%%%%%%%%%%%%%%%%%%%%%%%%%%%%%%%%%%%%%%%%

\section{Derivate parțiale}

Dată o funcție de mai multe variabile, se pot calcula derivatele parțiale ale
acesteia în funcție de fiecare dintre variabile, pe rînd. De exemplu,
fie funcția:
\[
  f : \RR^2 \to \RR, \quad f(x, y) = 3x \sin y + e^y + 4xy.
\]
Calculăm:
\[
  \frac{\partial f}{\partial x} = 3 \sin y + 4y, \quad %
  \frac{\partial f}{\partial y} = 3x\cos y + e^y + 4x.
\]

Mai departe, putem calcula și derivatele de ordin superior, după
regula:
\[
  \frac{\partial^2 f}{\partial x^2} = \frac{\partial}{\partial x} %
  \Big( \frac{\partial f}{\partial x} \Big)
\]
și similar pentru celelalte derivate parțiale, adică $ \dfrac{\partial^2 f}{\partial y^2} $,
$ \dfrac{\partial^2 f}{\partial x \partial y} $, respectiv
$ \dfrac{\partial^2 f}{\partial y \partial x} $.

De asemenea, pentru o funcție de două variabile $ f = f(x, y) $, se definește
\emph{laplacianul} funcției ca fiind:
\[
  \Delta f = \frac{\partial^2 f}{\partial x^2} + \frac{\partial^2 f}{\partial y^2}.
\]
Definiția se extinde natural pentru $ n $ variabile.

Dacă $ f $ este astfel încît $ \Delta f = 0 $, atunci funcția $ f $
se numește \emph{armonică}. \index{funcții!laplacian} \index{funcții!armonice}

În plus, dată o funcție de două variabile $ f = f(x, y) $, se definește
\emph{diferențiala totală} $ df $ ca fiind:
\[
  df = \frac{\partial f}{\partial x} dx + \frac{\partial f}{\partial y} dy.
\]

Similar se poate defini pentru funcții de $ n $ variabile, precum și
diferențiala totală de ordinul 2:
\[
  d^2 f = \frac{\partial^2 f}{\partial x^2} dx^2 + \frac{\partial^2 f}{\partial y^2} + %
  2 \frac{\partial^2 f}{\partial x \partial y} dxdy.
\]

%%%%%%%%%%%%%%%%%%%%%%%%%%%%%%%%%%%%%%%%%%%%%%%%%%%%%%%%%%%%%%%%%%%%%%

\section{Exerciții}

0. Calculați următoarele derivate parțiale folosind definiția:
\begin{enumerate}[(a)]
\item pentru $ f(x, y) = 3x^2 - 2xy + 5y^2 - 3x + 2 $, $ \dfrac{\partial f}{\partial x}(1, 2) $;
\item pentru $ f(x, y) = 5x^2 + xy - 3x + 2 $, $ \dfrac{\partial f}{\partial x}(-1, 1) $ și
  $ \dfrac{\partial f}{\partial y}(1, 2) $;
\item pentru $ f(x, y) = 3x - 2y^2 + 2xy + 5 $, $ \dfrac{\partial f}{\partial x}(0, 1) $ și
  $ \dfrac{\partial f}{\partial y}(-1, 0) $.
\end{enumerate}

1. Să se calculeze derivatele parțiale de ordinul întîi pentru funcțiile:
\begin{enumerate}[(a)]
\item $ f(x, y) = e^{x - y^2} $;
\item $ f(x, y) = \arctan \dfrac{x}{y} + \arctan \dfrac{y}{x} $;
\item $ f(x, y) = e^{x^2 + y^2} $;
\item $ f(x, y, z) = \sqrt{x^2  + y^2 + z^2} $;
\item $ f(x, y) = x^{\sqrt{y}} $.
\end{enumerate}

\vspace{1cm}

2. Calculați derivatele parțiale de ordinul întîi pentru funcțiile compuse:
\begin{enumerate}[(a)]
\item $ f(x, y) = \ln(u^2 + v), u(x, y) = e^{x + y^2}, v(x, y) = x^2 + y $;
\item $ f(x, y) = \varphi(2xe^y + 3y \sin 2x) $;
\item $ f(x, y) = \varphi(u, v, w) $, unde $ u(x, y) = xy, v(x, y) = x^2 + y^2 $ și
  $ w(x, y) = 2x + 3y $;
\item $ f(x, y) = \arctan \dfrac{2u}{v} $, unde $ u(x, y) = x\sin y $ și
  $ v(x, y) = x \cos y $;
\end{enumerate}

\vspace{1cm}

3. Să se calculeze derivatele parțiale de ordinul întîi și al doilea
pentru funcțiile:
\begin{enumerate}[(a)]
\item $ f(x, y) = e^x \cos y $;
\item $ f(x, y) = \dfrac{x - y}{x + y}, (x, y) \neq (0, 0) $;
\item $ f(x, y) = \ln(x^2 + y^2), (x, y) \neq (0, 0) $;
\item $ f(x, y) = x^3 + xy $;
\item $ f(x, y, z) = y \sin (x + z) $;
\item $ f(x, y, z) = \ln(x^2 + y^2 + z^2), (x, y, z) \neq (0, 0, 0) $.
\end{enumerate}

\vspace{1cm}

4. Să se calculeze derivatele parțiale în punctele indicate:
\begin{enumerate}[(a)]
\item $ f(x, y) = 2x^2 + xy, \dfrac{\partial f}{\partial x}(1, 1), %
  \dfrac{\partial f}{\partial y} (3, 2) $;
\item $ f(x, y) = \sqrt{\sin^2 x + \sin^2 y}, \dfrac{\partial f}{\partial x} %
  \Big( \dfrac{\pi}{4}, 0 \Big), \dfrac{\partial f}{\partial y} %
  \Big( \dfrac{\pi}{4}, \dfrac{\pi}{4} \Big) $;
\item $ f(x, y) = \ln(1 + x^2 + y^2), \dfrac{\partial f}{\partial x}(1, 1), %
  \dfrac{\partial f}{\partial y}(1, 1) $;
\item $ f(x, y) = \sqrt[3]{x^2 y}, \dfrac{\partial f}{\partial x}(-2, 2) $,
  $ \dfrac{\partial f}{\partial y}(-2, 2) $, $\dfrac{\partial^2 f}{\partial x\partial y}(-2, 2) $;
\item $ f(x, y) = xy\ln x, x \neq 0, \dfrac{\partial^2 f}{\partial x\partial y}(1, 1) $,
  $ \dfrac{\partial^2 f}{\partial y}{\partial x}(1, 1) $.
\end{enumerate}

\vspace{1cm}

5. Arătați că funcțiile următoare verifică ecuațiile indicate:
\begin{enumerate}[(a)]
\item $ f(x, y) = \varphi \Big( \dfrac{y}{x} \Big) $,
  \[
    x \frac{\partial f}{\partial x} + y \frac{\partial f}{\partial y} = 0.
  \]
\item $ f(x, y, z) = \varphi(xy, x^2 + y^2 + z^2) $,
  \[
    xz \frac{\partial f}{\partial x} - yz \frac{\partial f}{\partial y} + %
    (y^2 - x^2) \frac{\partial f}{\partial z} = 0.
  \]
\item $ f(x, y) = y \varphi(x^2 - y^2) $,
  \[
    \frac{1}{x} \frac{\partial f}{\partial x} + \frac{1}{y} \frac{\partial f}{\partial y} = %
    \frac{1}{y^2} f(x, y);
  \]
\item $ f(x, y, z) = \dfrac{xy}{z} \ln x + x\varphi \Big( \dfrac{x}{y}, \dfrac{z}{x} \Big) $,
  pentru $ x > 0 $ și $ z \neq 0 $, ecuația fiind:
  \[
    x \frac{\partial f}{\partial x} + y \frac{\partial y}{\partial y} + z \frac{\partial f}{\partial z} - %
    \frac{xy}{z} - f(x, y, z) = 0.
  \]
\end{enumerate}

\vspace{1cm}

6. Verificați dacă următoarele funcții sînt armonice:
\begin{enumerate}[(a)]
\item $ f(x, y) = \ln(x^2 + y^2) $;
\item $ f(x, y, z) = \ln(x^2 + y^2 + z^2) $;
\item $ f(x, y) = \sqrt{x^2 + y^2} $;
\item $ f(x, y, z) = \sqrt{x^2 + y^2 + z^2} $;
\item $ f(x, y) = (x^2 + y^2)^{-1} $;
\item $ f(x, y, z) = (x^2 + y^2 + z^2)^{-1} $.
\end{enumerate}


\vspace{1cm}

7. Reprezentați grafic domeniile date de următoarele (in)ecuații:
\begin{enumerate}[(a)]
\item $ D_1 = \{ (x, y) \in \RR^2 \mid y = x^2, y^2 = x \} $;
\item $ D_2 = \{ (x, y) \in \RR^2 \mid x = 2, y = x, xy = 1 \} $;
\item $ D_3 = \{ (x, y) \in \RR^2 \mid x^2 + y^2 \leq 2y \} $;
\item $ D_4 = \{ (x, y) \in \RR^2 \mid x^2 + y^2 \leq x, y \geq 0 \} $;
\item $ D_5 = \{ (x, y) \in \RR^2 \mid y = 0, x + y - 6 = 0, y^2 = 8x \} $;
\item $ D_6 = \{ (x, y) \in \RR^2 \mid x^2 + y^2 \leq 2x + 2y - 1 \} $.
\end{enumerate}



%%% Local Variables:
%%% mode: latex
%%% TeX-master: "../bcd-18-19"
%%% End:
